% SOURCE docs/release/2026-10-29/source/docs/writing-repair/2026-10-29/src/d97a1b72be46db93.tex CHAR 0
\chapter{Using the analysis in other research applications}
\label{in:consumers}
\label{exec13:i3-start}

A useful connection joins a consumer's decision to a provider's scientific result without changing either. The cases here follow four banking journeys: funding value, sequential servicing, discovery ambiguity and trustworthy observations. Population transport and cash-flow composition then show why correct individual exchanges may still fail when combined. The resulting comparison favors a reviewed manual exchange as the baseline and automation only where the same scientific output can be preserved at a useful total cost.


% SOURCE docs/release/2026-10-29/source/docs/writing-repair/2026-10-29/src/5598898bc32afb70.tex CHAR 0
\section{What a complete connection must preserve}

A treasury analyst may obtain a state estimate from one package, an intervention
contrast from another and a value calculation from a third. The final number
answers the original question only if these exchanges preserve its meaning.
An exposure measured in millions, a rate expressed in basis points, an outcome
available after a decision and a value computed under a different continuation
policy can each change that meaning while leaving a syntactically valid table.
The integration problem therefore begins with the banking decision and ends
with the consumer's displayed result. Successful transfer of a numeric array
is one intermediate condition.

Let the target declaration be
\[
 \mathcal T=(\mathcal P,\pi,\pi_0,t_0,H,\mathcal F,u,c),
\]
where $\mathcal P$ is the population, $\pi$ and $\pi_0$ are the intervention
and comparator, $t_0$ is the decision origin, $H$ the outcome horizon,
$\mathcal F$ the available information, $u$ the measurement unit and $c$ the
cash-flow convention. This is a proposed semantic record, rather than a claim
that every inspected repository already serializes these fields. A target also
needs its scientific premises: the observation and assignment laws, graph or
model class, sampling units, and the evidence supporting any restrictions.
Those premises are referenced alongside $\mathcal T$; a short tuple is not a
replacement for their mathematical definition.

Write $D$ for a consumer's data and $M(D)$ for the representation sent to the
causal operation. A correct adapter must give a declared mapping $g$ for which
\begin{equation}
 \tau_{\mathrm{consumer}}(D;\mathcal T)
   =g\{\tau_{\mathrm{provider}}(M(D);\mathcal T')\}
 \label{exec13:i3-target-map}
\end{equation}
under the stated model. For a unit conversion, $g$ may be a fixed affine
function. For a population change, time aggregation or policy substitution,
there is usually no such automatic identity. The adapter must then require a
new analysis or an additional argument. This requirement is useful even when
both packages use identical field names: an ``effect'' can denote an
association, a structural response or a policy-value difference.

The result mapping must preserve uncertainty as well as the target. For a
fixed known function $g$ and a confidence set $C$ for $\tau$, the event
$\{\tau\in C\}$ implies $\{g(\tau)\in g(C)\}$. Thus applying $g$ to the
whole set preserves the coverage lower bound. This elementary implication
does not make a supplied interval a valid confidence set, and it does not
apply to unaccounted uncertain exchange rates, estimated exposure weights or
a transformation chosen after inspecting the result. A joint confidence
region for all uncertain inputs, or an explicitly conditional scenario, is
needed in those cases. An identified set, a sensitivity range and a confidence
set must keep different result types even when each happens to be printed as
two endpoints.

Three connections consequently need separate judgments. Source inspection can
establish that a field or command is declared. An author-operated trial can
establish that a selected version transports and interprets a fixture
correctly. Acceptance by a consumer or an institution requires its own
decision and evidence. The cases below use inspected interfaces, existing
local runtime evidence and the transparent calculations developed in this chapter at their respective levels. The fitted studies elsewhere in the book provide additional evidence for particular statistical regimes, but do not qualify these proposed consumer mappings. They do not claim that sibling maintainers have adopted this library.

\section{The workspace as a set of scientific roles}

The dated inventory observed eighteen top-level directories and inspected
fifty-three source records on 12 October 2026. It recorded three MacroFinance
directories as one project family with distinct revisions. Its observation
date is later than the 7 October literature cutoff, so it establishes a local
environment snapshot rather than a retrospective software census~\citep{v12_in_inventory}. Selected declarations for FinMarketing, CardNPV,
BayesFilter and dataml were inspected again on 13 October. A comparison of all
fifty-three historical identities found ten changed files; the four selected
technical declaration files retained their recorded bytes. Historical
descriptions of the other projects therefore remain dated observations, not
claims about their latest behavior~\citep{exec13_i3_inspection}.

Table~\ref{exec13:i3-workspace} gives a disposition for every observed
directory. The distinctions are architectural. A numerical provider supplies
a likelihood, state or solver result; a consumer supplies a question and uses
the answer; an evidence tool locates or checks material. Mixing these roles
would let a solver's convergence stand in for identification, or let a document
summary stand in for a primary source. The default location for new mappings
is the causal project, where their scientific restrictions can be reviewed
together. No cross-repository import is required merely to describe a proposed
exchange.

\begin{longtable}{P{.21\linewidth}P{.23\linewidth}P{.47\linewidth}}
\caption{Dated workspace dispositions and the proposed next connection.
All interfaces require version-specific confirmation before use.}
\label{exec13:i3-workspace}\\
\toprule Directory & Role and choice & Reason and next scientific boundary \\
\midrule\endfirsthead
\toprule Directory & Role and choice & Reason and next scientific boundary \\
\midrule\endhead
BayesFilter & Numerical provider; integrate selected outputs & Retain model,
data, coordinates, approximation and diagnostics; require a separate
intervention argument. \\
MacroFinance & Funding consumer; select one model & Use its designated
BayesFilter authority, with action basis, horizon, exposure and structural
meaning declared. \\
MacroFinance worktrees & Group, retain revisions & The directories
\path{MacroFinance-ccma-g-v7-neutra-r9} and
\path{MacroFinance-dz5-neutra} are separate checkouts within this family;
neither is an additional independent product comparison. \\
FinMarketing & Priority consumer & Map source, feature, target and value
records while preserving source ownership, availability and decision clocks. \\
dataml & Priority observation exchange & Review artifact and assumption
records before translating them; declared JSON output is a transport boundary,
not verified causal truth. \\
cardnpv & Valuation provider or consumer & Preserve historical minus
perfect-execution value, admitted horizon and continuation; do not rename it
as a policy effect. \\
DynareMCP & Model provider; conditional integration & Retain model equations
and observation timing; its inspected CLI/in-process profile does not establish
strict stdio conformance. \\
MathDevMCP & Mathematical evidence tool & Pass explicit obligations and
premises; inspect its restricted-use terms and transport before adoption. \\
research-assistant & Source navigation & Consider the inspected read-only
route; retain genuine PDFs and technical anchors outside generated summaries. \\
SMEwallet & Relevant; interface deferred & The inspected mathematical
primitives are not an estimator. Specify its target/accounting bridge before
selecting an adapter. \\
MOONeural & Adjacent solver; defer & Require a concrete nuisance or numerical
task and applicable rights; a provenance notice alone is no license grant. \\
hponas & Adjacent tuning; defer & Require a declared learner, training split
and license review; tuning success does not qualify causal inference. \\
humanvoice & Editorial; defer & A later prose revision must preserve claims
and citations; editorial changes cannot provide scientific acceptance. \\
dsge\_hmc & Container; defer & Three nested checkouts were observed. Select a
canonical model and interface before treating the container as a consumer. \\
claudecodex & Exclude from causal runtime & Shared agent policy infrastructure
supplies no causal algorithm or bank observation interface. \\
.local & Exclude support directory & Ambient environment files are not a
declared application dependency; contents were outside the inspection. \\
causal & Producer and standalone client & Preserve the distinction between
program-control commands and the scientific library, CLI and MCP surfaces. \\
\bottomrule
\end{longtable}

The initial consumer shortlist is FinMarketing, dataml and MacroFinance,
with BayesFilter and CardNPV as additional numerical and value boundaries.
This is a proposal about useful first connections. It does not imply that all
five should be loaded into one process. A file exchange is often the clearest
first design: it keeps the original target, source identity and errors visible
without coupling environments. A subprocess can later automate a selected
exchange while preserving an isolated environment. A direct library adapter
becomes attractive when repeated serialization dominates a measured workflow
and the object semantics have already been demonstrated. Porting a numerical
implementation creates additional maintenance and scientific equivalence
obligations; it is not justified merely by a common programming language.

License and ownership conditions belong to this choice. The inspected
CardNPV and MathDevMCP notices described internal-use restrictions; other
entries had incomplete or proprietary declarations~\citep{v12_in_inventory}.
The comparison supports designing exchanges, not a distribution grant.
Where terms are unresolved, a documented manual export remains a design
alternative, subject to the exporter having authority to supply the material.
A prospective port must separately identify the code to be copied, its
license obligations and the cost of maintaining its scientific behavior.

\section{Information clocks before feature exchange}

A transaction can occur before a decision and still be unavailable to the
decision maker. Let $v_j$ be the event time for feature $j$, $a_j$ its
availability time and $d_t$ the action time. The relevant information set is
generated by records available before $d_t$, not by every record whose event
time precedes it. Late corrections, external bureau arrivals and retrospective
entity resolution all create this distinction. Reconstructing the final
database as though it were visible at every earlier date gives a different
policy from the one the bank could have used.

FinMarketing's inspected \texttt{FeatureSnapshot} compares timestamp and
sequence pairs and requires
\[
 (a_{\max},s_a)\leq(c,s_c)<(d_t,s_d),
\]
where $c$ is the feature cutoff and $s$ is a within-timestamp sequence.
Its selected records also distinguish training, validation, test and serving
partitions, bind a fitted transformation to a snapshot, and keep outcome
names separate from feature names. The inspected \texttt{SourceBatch}
includes owner, use basis, permission and source digest, with a declared
synthetic or licensed-public boundary~\citep{exec13_i3_inspection}. These
declarations give a concrete proposed mapping. They do not establish that
the supplied timestamps or ownership claims are true.

Consider an illustrative decision pair $(100,4)$ and cutoff $(100,3)$.
Availability at $(99,9)$ or $(100,3)$ is admitted; availability at $(100,5)$
or $(101,0)$ is rejected. The four exact checks are reproduced with the other
chapter examples~\citep{exec13_i3_examples}. Dropping the sequence field makes
the third case appear simultaneous and loses the needed ordering. Converting
UTC to a local display time must not change the comparison. When a source
does not supply the required ordering, the adapter should request a coarser
defensible decision clock or report unavailable information, rather than
inventing the missing sequence.

Filtering introduces the same issue at a less visible boundary. A state
estimate using observations through $t$ can support a decision after those
observations arrive. A smoothed state using later observations cannot be
substituted into that historical policy without changing its information
set. Both arrays can have the same dimensions and plausible values. The
consumer therefore needs the conditioning horizon and observation-availability
rule, in addition to a state name and numeric uncertainty. This requirement
connects the workspace interfaces directly to the longitudinal assumptions
developed earlier in the monograph.


% SOURCE docs/release/2026-10-29/source/docs/writing-repair/2026-10-29/src/fff981099c670a3e.tex CHAR 0
\section{Funding: from a structural rate response to value}

Suppose treasury asks for the value of a ten-basis-point reduction in a
funding action over four quarters. Three objects are required. The consumer
must specify the action and exposure. A scientific analysis must establish
what change in paid funding rate that action would cause under its model.
A valuation step must then translate that response into dated cash flows.
A forecast of market rates alone supplies neither the intervention response
nor the exposure response to the bank's decision.

The inspected MacroFinance inventory directs active numerical work to
BayesFilter and describes the root migration adapter as a compatibility
facade~\citep{v12_in_inventory}. In the re-inspected BayesFilter declarations,
\texttt{BayesianSSMProblem} binds the model manifest, data signature, static
shape and coordinate convention; \texttt{FilterProgram} separately names the
filter, required model capabilities, target policy and approximation
semantics~\citep{exec13_i3_inspection}. A proposed funding adapter should retain
these distinctions. An exact structural model, a deterministic approximate
filter and a fixed-randomness numerical procedure need not imply the same
likelihood or uncertainty. Importing a historical local sampler because it
returns a familiar array could bypass the selected authority entirely.

For a transparent conditional example, fix exposure $B=100$ million USD,
action $a=-10$ basis points, annual effective discount rate $r=0.04$ and an
initial implementation cost $K=40{,}000$ USD. Let $\theta$ denote paid-rate
pass-through in basis points per action basis point, held constant across
four quarters. Assume balances are invariant to the action, cash accrues
uniformly within each quarter and is paid at quarter end, and no additional
fees or risk effects occur. The present value of savings is then
\begin{equation}
 V(\theta)=-B\frac{a}{10{,}000}\theta
      \sum_{q=1}^{4}\frac{1/4}{(1+r)^{q/4}}-K
       =97{,}584.369097\theta-40{,}000.
 \label{exec13:i3-funding}
\end{equation}
The sign is explicit: a negative rate change decreases cost and produces a
positive saving. The factor $10{,}000$ converts basis points to a decimal
rate, while $1/4$ converts an annualized rate into quarterly accrual. The
fractional discount exponent is required because $r$ is annual effective.
Replacing it with a quarterly rate would require a different parameterization.

For an assumed pass-through range $[0.6,0.9]$, the resulting value range is
$[18{,}550.62,47{,}825.93]$ USD; at $\theta=0.75$ the value is
$33{,}188.28$. The break-even coefficient is approximately $0.409902$.
These figures reproduce both by the affine formula and by summing the four
cash flows, including a dollars-to-millions round trip. They are generated
planning calculations with supplied inputs, not bank estimates or a
confidence interval~\citep{exec13_i3_examples}. The coefficient range could
instead be an identified set or a fitted confidence set, but only if a
scientific analysis justified and labeled that interpretation.

The example makes several adverse exchanges visible. Treating each quarter
as a full year of accrual changes the midpoint value to $252{,}753.11$.
Reading ten basis points as ten percentage points changes the rate scale by
a factor of one hundred. Reading millions as dollars changes the exposure
scale by a factor of one million. These are semantic errors even if the
downstream computation is numerically exact. A consumer acceptance fixture
should assert the declared dimensions and the independently computed cash
flows, not only compare two calls to the same valuation function.

More consequentially, an action may change the balance itself. With
action-dependent balance $B_q(a)$ and funding rate $s_q(a)$, the cost
contrast is based on
\[
 \sum_q d_q\Delta_q
       \{B_q(a)s_q(a)-B_q(0)s_q(0)\},
\]
where $d_q$ discounts the cash flow and $\Delta_q$ is the accrual fraction.
Equation~\eqref{exec13:i3-funding} follows only under its fixed-exposure
restriction and the specified rate contrast. Multiplying a rate effect by
the observed balance does not account for endogenous deposit attrition,
replacement funding or interaction between amount and price. A useful
adapter therefore either retains this restricted target or asks for the
joint structural response. It does not silently convert a restricted
scenario into a general treasury value claim.

The first proposed workflow is consequently modest and complete: a
MacroFinance question supplies target and exposure declarations; a selected
numerical output supplies its model and approximation identities; a causal
analysis accepts or refuses the intervention meaning; and a value object
retains the transformation and the complete uncertainty object. A simpler
baseline is a manually reviewed exported result with the same declarations.
Only a measured reduction in repeated translation errors or effort, with
equivalent scientific output, supports replacing that baseline with an
automated connection. No general MacroFinance--BayesFilter--causal adapter
has been qualified by the arithmetic example.

\section{Numerical authority and the error in a value transformation}

A connection to a state-space provider also needs a mathematical account of
approximation. Suppose, conditional on an explicitly declared model and
available history, the provider's target distribution over a finite state
$S$ is $P$ and its numerical approximation is $\widetilde P$. Let $f(S)$
be the value contribution, with known bounds $m\leq f(s)\leq M$.
Define total variation by
\[
 \operatorname{TV}(P,\widetilde P)
     =\frac12\sum_s|P(s)-\widetilde P(s)|.
\]
Then a supplied bound $\operatorname{TV}(P,\widetilde P)\leq\epsilon$
implies
\begin{equation}
 |E_Pf-E_{\widetilde P}f|\leq(M-m)\epsilon.
 \label{exec13:i3-tv}
\end{equation}
For a direct proof, let $\delta_s=P(s)-\widetilde P(s)$. The sum of the
positive $\delta_s$ and minus the sum of the negative $\delta_s$ both
equal total variation, since $\sum_s\delta_s=0$. Subtract the constant
$m$ from $f$, and bound the positive contribution by
$(M-m)\sum_{\delta_s>0}\delta_s$. Interchanging $P$ and $\widetilde P$
gives the other direction. This bound is a project derivation for the
specified finite problem, not a claim that an inspected filter supplies
$\epsilon$ automatically.

In the funding example, let the low and high states have coefficients
$0.6$ and $0.9$. Suppose the declared target probability of the high state
is $0.4$, while the approximate provider uses $0.5$. Their mean
coefficients are $0.72$ and $0.75$, and total variation is $0.1$.
The target conditional value is $30{,}260.75$ USD and the approximate
value is $33{,}188.28$ USD. Their difference $2{,}927.53$ equals the
bound in Equation~\eqref{exec13:i3-tv}, because the value range is
$29{,}275.31$. Exact rational coefficient arithmetic and a common-uniform
coupling check the calculation independently~\citep{exec13_i3_approximation}.
Neither distribution here is output from an executed BayesFilter run.

This construction shows what the consumer would need. If only the
approximate value and a justified $\epsilon=0.1$ bound were supplied,
the corresponding symmetric outer value interval would be
$[30{,}260.75,36{,}115.81]$ USD. That interval concerns numerical error
relative to the declared target distribution. It is not a frequentist
confidence interval for a bank effect and is not a posterior credible
interval for the random state. Model uncertainty, parameter uncertainty
and identification ambiguity require their own treatment. If the payoff is
unbounded, the displayed finite-range proof is unavailable; truncation,
tail control or another error metric would have to be justified rather
than assumed.

For an actual exchange, the causal-side mapping should therefore retain
the following fields in distinct roles. The model/data/coordinate identity
specifies which conditional distribution is being approximated. The
filter identity and approximation semantics identify the numerical
procedure. Its diagnostics describe what was observed in that procedure,
while an error statement must specify a norm, target and validity regime.
Finally, the value mapping gives the payoff, its horizon and any bounds
used to translate numerical error. A convergence trace is useful diagnostic
evidence but cannot silently become a total-variation certificate. When no
justified error bound is available, the consumer should display a
conditional approximate result and the missing error qualification.

A similar distinction applies to posterior output. A posterior distribution
computed under a structural model may support Bayesian decision analysis
conditional on that model and its priors. Passing posterior quantiles into
a field labeled frequentist coverage changes the asserted guarantee. Passing
a likelihood under a merely predictive model into an intervention-value
operation changes the scientific interpretation even earlier. The proposed
common contract should keep these tags and premises with the object, so a
consumer can compare a conditional structural calculation with a design-based
alternative without presenting them as interchangeable evidence.

This suggests a useful progression for the first numerical-provider trial.
Start with a finite model whose state law and value are exactly enumerable.
Check the declared approximation and the monetary mapping against that
reference. Next vary numerical settings while keeping the target fixed,
recording error, memory and time separately. Only then consider a richer
model with an explicit replacement for the finite reference. Such a sequence
tests the connection's mathematical meaning at low cost. It leaves the
intervention assumptions and the usefulness of its precision visible for a
later substantive decision.

\section{Servicing: preserve both policies and both decision times}

A servicing question asks how a feasible sequence of actions changes value.
It cannot be answered merely by importing an account's observed outcome or
a retrospective value gap. To make the required exchange concrete, consider
the following generated two-action law. The binary baseline covariate $X$
has probability $1/2$ of either value. Action $A_1$ occurs first; the binary
intermediate response $L$ then becomes available before action $A_2$; the
binary final recovery indicator $Y$ arrives last. Under interventions on
the two actions, independent exogenous uniforms generate
\begin{align}
 \Pr(L=1\mid X=x,A_1=a_1)&=0.2+0.2x+0.2a_1,\\
 \Pr(Y=1\mid x,a_1,l,a_2)&=0.1+0.2x+0.3l+0.2a_2+0.1a_1.
 \label{exec13:i3-servicing-law}
\end{align}
All stated probabilities lie between zero and one. Recovery earns 100 USD;
the first and second actions cost 6 and 4 USD respectively. For this short
illustration discounting is zero, and the total reward is
$R=100Y-6A_1-4A_2$. The independence and intervention interpretation are
definitions of this generated structural law, not conclusions from a fitted
observational regression.

Compare the policy $\pi_0$ taking neither action with the adaptive policy
\[
 \pi_1(x)=1-x,\qquad \pi_2(x,a_1,l)=1-l.
\]
The first action can use $X$ but not the later $L$. The second action may
use $L$ because its availability is part of the design. Write $p_L$ and
$p_Y$ for the two displayed probabilities. Direct integration gives
\begin{equation}
 V(\pi)=\frac12\sum_{x=0}^{1}\sum_{l=0}^{1}
   p_L(x,\pi_1(x))^l\{1-p_L(x,\pi_1(x))\}^{1-l}
   \left[100p_Y(x,\pi_1(x),l,\pi_2)-6\pi_1(x)-4\pi_2\right],
 \label{exec13:i3-policy-value}
\end{equation}
with $\pi_2$ evaluated at the displayed history. The sum integrates over
the intermediate response distribution induced by the first action. Holding
the observed distribution of $L$ fixed while changing $A_1$ would calculate
a different object. Substitution yields $V(\pi_0)=29$ and
$V(\pi)=43.6$, so the proposed contrast is $14.6$ USD per account.

An independent enumeration checks the same quantity through an observed
behavior law. Let $\Pr(A_1=1\mid X)=1/2$ and let
$\Pr(A_2=1\mid X,A_1,L)$ equal $1/4$ when $L=1$ and $3/4$ when
$L=0$. Generate actions with additional independent uniforms, so sequential
exchangeability and positivity hold by construction. If $e_1,e_2$ denote
the probabilities of the actions actually taken, then
\begin{equation}
 V(\pi)=E\left[
  \frac{\mathbf1\{A_1=\pi_1(X),\ A_2=\pi_2(X,A_1,L)\}}
       {e_1(A_1\mid X)e_2(A_2\mid X,A_1,L)}R\right].
 \label{exec13:i3-policy-ipw}
\end{equation}
To see the identity, factor the joint law of $(X,A_1,L,A_2,Y)$ and cancel
the two action probabilities on histories matching the policy. Summing over
$Y$ then gives its conditional mean, leaving
Equation~\eqref{exec13:i3-policy-value}. Exact rational enumeration of all
32 observed cells verifies total mass one, both values and the contrast~\citep{exec13_i3_examples}. This is a check of the generated finite law and
the proposed semantics. It supplies neither empirical coverage nor an
implemented longitudinal consumer adapter.

The corresponding FinMarketing exchange needs two feature snapshots or an
equivalent explicit history representation, two action definitions, the
availability of $L$, the cash horizon and the training identity of any fitted
policy. A policy chosen using the final evaluation outcomes is a further
target-selection problem, even when every field passes a schema. Candidate
policy learning should use an independent training sample or a separately
justified evaluation design. A filtering provider can assist in constructing
a declared observable history; it cannot make a privately observed state
available to the bank or remove unmeasured confounding by returning a latent
state estimate.

CardNPV illustrates a different target boundary. Its inspected
\texttt{CardTargetValue} uses historical value minus perfect-execution
value, supports the declared horizons of 24, 48, 120 and 240 months, and
distinguishes finite-only from conditional continuation~\citep{exec13_i3_inspection}.
Suppose, at a common admitted horizon, historical value is 900 USD and
perfect-execution value is 1,100 USD. That target is $-200$ USD. If a
particular feasible policy has value 980 USD, its contrast against history
is $80$ USD. Negating the accounting gap gives $200$, which still refers
to the perfect-execution intervention rather than the feasible policy.
An adapter must therefore compare the interventions, populations, horizons
and continuation assumptions before relating the targets. The generated
two-action example has its own short horizon; it is not supplied to CardNPV
as though it were an admitted 24-month account calculation.

Three alternatives remain useful. A scenario valuation can hold a supplied
policy response fixed and clearly label its assumptions. A longitudinal
causal analysis can attempt to identify the response from declared data and
assignment conditions. A structural state model can supply a conditional
intervention calculation under stronger dynamics. These alternatives answer
different evidential questions. The proposal should compare their assumptions,
information cost and diagnostics before selecting one for an actual consumer.
If the local service supports only a static target, the correct response to
this sequential request is an explicit unsupported-operation result together
with the preserved request; it must not return a static ATE under the same
policy label.


% SOURCE docs/release/2026-10-29/source/docs/writing-repair/2026-10-29/src/53f2135cc3868665.tex CHAR 0
\section{Discovery: carry ambiguity through to the financial decision}

A model or discovery provider can return candidate diagrams without
identifying the effect requested by a business user. The proposed
DynareMCP exchange would retain variable definitions, equation and observation
timing, and distinctions between structural and measurement variables. The
inspected documentation supported selected CLI/in-process operation while
withholding strict stdio qualification~\citep{v12_in_inventory}. Thus even the
transport requires a selected implementation and trial. More fundamentally,
neither the presence of structural equations nor the name of a discovery
method proves that an imported edge has an intervention interpretation.

A two-variable example isolates the issue. Let binary variables $A,Y$ have
joint probabilities
\[
 \begin{array}{c|cc}
       &Y=0&Y=1\\\hline
 A=0&0.4&0.1\\
 A=1&0.1&0.4
 \end{array}.
\]
One causal model generates $A$ fairly and then generates $Y$ with probability
$0.2+0.6A$, using an independent uniform disturbance. Its graph is
$A\to Y$, and $E\{Y\mid\operatorname{do}(A=1)\}-
E\{Y\mid\operatorname{do}(A=0)\}=0.6$. A second model generates $Y$
fairly and then $A$ with probability $0.2+0.6Y$. Its graph is $Y\to A$,
and intervening on $A$ changes no equation for $Y$, giving effect zero.
Multiplication of the two factors verifies exactly the same joint table
in both models~\citep{exec13_i3_examples}.

In the absence of additional information distinguishing these two models,
the proposed consumer must preserve both effects. If one additional recovery
is valued at 100 USD and an intervention costs 20 USD per person, the two
net values are $40$ and $-20$ USD. Their signs differ. Choosing the arrow
that produces a favorable investment case would therefore answer a different,
more restricted question. The set $\{0,0.6\}$ describes these two specified
models; it is not asserted to exhaust all latent-variable models or to be a
sampling confidence region. Its convex hull $[0,0.6]$ is a conservative
outer summary, but loses the fact that interior effects do not occur in this
particular two-model calculation.

This example is deliberately free of sampling error. More observations from
the same joint law cannot distinguish the two constructions. A real time
ordering, an intervention, or a justified scientific restriction can change
the admissible class. For example, a documented action that genuinely precedes
the outcome would rule out this literal reverse-time model, although it
would not by itself exclude hidden common causes. A language model's preferred
wording, a graph rendering package or a high prediction score supplies none
of those restrictions. Evidence used to remove a graph must be retained with
the result and with the specific variables and time indices it concerns.

The discovery-to-effect exchange should consequently carry a class object:
node identities and measurement definitions, endpoint semantics, permitted
latent variables and time lags, restrictions with sources, and the exact
analysis target. A causal operation can return effects across admitted graphs
when its identification domain supports them. Otherwise it should return the
supported partial result and a reason for withholding the effect. A consumer
that accepts only a single DAG needs an explicit narrowing decision; selecting
one representative for display must not narrow the scientific object passed
to estimation. The Java and R adapter discussions earlier in the book develop
this same requirement for specific graph representations.

MathDevMCP has a possible supporting role at the derivation boundary: a
client could send a located identity, its assumptions and a requested
verification operation, then retain the returned evidence. The inspected
project distinguishes specialist verification from generated derivations and
has restricted-use terms~\citep{v12_in_inventory}. A computer algebra check of
the two factorizations can confirm their arithmetic. It cannot certify that
a bank's actual data were generated by either model. The simplest exact
enumeration is sufficient for the present counterexample; a general
mathematical service would need to justify its additional complexity.

\section{Trustworthy observations: selection and origin identity}

A clean table may preserve measurements while changing the population
represented. The proposed dataml--FinMarketing exchange therefore begins
with origin identity, source owner, inclusion mechanism, record disposition,
availability and any unresolved linkage alternatives. The inspected dataml
research CLI declares JSON on standard output, nonzero failures, and separate
call, export and MCP entrypoints~\citep{exec13_i3_inspection}. These declarations
are useful for arranging an export contract, but a JSON success response
cannot establish that an inclusion probability is correct or that a label is
truthful. An adapter must preserve the source of those declarations and the
scope of their verification.

Consider a finite population of four origins with values
$(20,40,60,80)$ and independent inclusion indicators $R_i$ with probabilities
$(1/2,1,1/4,1/2)$. The target is the equally weighted population mean 50.
If origin identities and probabilities are known and labels are correct,
the expansion estimator
\begin{equation}
 \widehat\mu=\frac14\sum_{i=1}^4\frac{R_iY_i}{p_i}
 \label{exec13:i3-selection}
\end{equation}
has expectation $\frac14\sum_iY_i=50$, since $E(R_i/p_i)=1$.
Under the specified independent selection law,
\[
 \operatorname{Var}(\widehat\mu)
    =\frac1{16}\sum_i\frac{1-p_i}{p_i}Y_i^2=1100.
\]
This exact variance shows that unbiasedness alone does not make the small
sample precise. It is a design variance for a fixed finite population, not
an estimated variance from four independent new bank customers.

Enumerating all sixteen possible inclusion vectors verifies both moments.
The second origin is always included, so the unweighted observed mean is
defined whenever the design has positive probability. Its expectation is
$365/8=45.625$, rather than 50~\citep{exec13_i3_examples}. Simply discarding
the selection fields and treating the exported table as representative
therefore changes the answer. A source may instead intentionally ask for the
mean among selected records. That is a legitimate different target, which
should carry a different population declaration.

If the third origin has inclusion probability zero, the problem changes
again. The populations $(20,40,60,80)$ and $(20,40,100,80)$ then give the
same complete observed-data law under the same selection of the other
origins, but their means are 50 and 60. No weighting formula can use a
never-observed origin to distinguish these worlds. Additional outcome bounds,
a changed sampling design or new information are necessary. A numeric
substitute for division by zero would conceal this identification gap.
Similarly, estimating the $p_i$ creates another uncertainty source, and
correlated selections require joint inclusion information for the variance.
Equation~\eqref{exec13:i3-selection} does not resolve either issue merely
because a probability field is present.

Linkage has a related but separate dependence problem. Suppose two monetary
errors each take values $-10$ and $10$ with equal marginal probability. If
they are the same shared mapping error, the variance of their sum is 400;
if they are independent, it is 200. The marginal correctness description is
identical. An exchange that gives each record its own independent error bar
would discard consequential shared uncertainty. The adapter should retain
common origin, household, calendar and mapping identifiers when the selected
analysis requires them. A common identifier indicates a possible dependence
structure; it is not itself an estimated covariance or a proof of an
independence assumption.

For the proposed trustworthy-observation journey, the consumer first states
the target population and information clock. The source provider then exports
records, inclusion and label declarations, unresolved ownership or matching
alternatives, and accepted/rejected counts. The causal layer checks whether
its method supports that information structure and returns either an
analysis with the retained limitations or a precise missing-information
result. Its displayed output must show when the conclusion is conditional
on supplied selection or truth assertions. Authentic ownership and bank gold
labels remain requirements for later claims about a bank; the present
proposal can demonstrate these logical distinctions using public or generated
information without pretending to have obtained that validation.

Research-assistant can help locate documentation and original papers for
these declarations through its inspected read-only source-navigation route~\citep{v12_in_inventory}. A source summary should point back to a genuine local
document and the inspected technical passage. It cannot replace the paper,
prove a data owner's assertion or authorize a new use of the records. This
separates a useful navigation integration from the much stronger and
unsupported claim of automated scientific truth reconstruction.


% SOURCE docs/release/2026-10-29/source/docs/writing-repair/2026-10-29/src/41c7e5b763063281.tex CHAR 0
\section{Population changes are new analyses}

Two consumers can agree on an action and outcome while referring to different
customers. A marketing sample might overrepresent a low-risk stratum, while a
servicing portfolio contains mostly a high-risk stratum. Changing the
population field in a copied result would not transport the effect. The
inspected FinMarketing comparison requires equality of population, action,
comparator, origin, horizon, finance owner and capital convention when mapping
an estimate into value; its comment explicitly distinguishes a unit conversion
from reaggregation across populations or horizons~\citep{exec13_i3_inspection}.
This gives the proposed adapter a useful boundary: a new population needs
a new target identity and a scientific bridge.

For a simple bridge, let $X$ have two strata, let $P$ describe the source
population and $Q$ the target population, and suppose the stratum-specific
causal effect $\tau(x)$ is invariant between them. This is an explicit
transport assumption in addition to identification of the source effects.
If $Q(x)>0$ only where $P(x)>0$, then
\begin{equation}
 \tau_Q=\sum_xQ(x)\tau(x)
       =\sum_xP(x)\frac{Q(x)}{P(x)}\tau(x).
 \label{exec13:i3-transport}
\end{equation}
The equality is algebraic under the stated assumptions. It does not establish
that equal stratum names represent equal measurements, that treatment versions
match, or that effects are in fact invariant. Different credit limits,
collection actions or outcome definitions would require a richer bridge.

Suppose the stratum effects are $10$ and $-10$ USD per account. With source
weights $(0.8,0.2)$ the source average is $6$ USD. With target weights
$(0.2,0.8)$ the target average is $-6$ USD. Density ratios $(1/4,4)$
reproduce the second value from the source distribution exactly~\citep{exec13_i3_composition}. The sign reversal arises entirely from the
population change, despite identical conditional effects. A common API and
identical units would not prevent it. If the second stratum is absent from
the source, the target result is not obtained by capping the missing ratio.
One must restrict the target, supply additional information or retain a bound
under stated restrictions.

For a FinMarketing or SMEwallet consumer, a proposed export should therefore
identify the unit of analysis, stratum construction, population counts or
weights, intervention versions and availability of the variables used to
transport. If estimated effects and estimated population weights are combined,
their joint uncertainty matters. Splitting the two estimates across
repositories does not make their errors independent. A consumer can instead
request a conditional planning calculation with fixed supplied weights, but
the response should preserve that qualification. The public-data and transport
chapters discuss alternatives when the desired bank population is not yet
observed; a bank-specific transported claim is not needed to develop this
proposal.

\section{Portfolio composition and the ownership of cash flows}

Separate funding and servicing effects also need not add. Consider a
generated net-value table for one common account, with funding action $F$
and servicing action $S$:
\[
 Y(f,s)=100+20f+15s-30fs,\qquad f,s\in\{0,1\}.
\]
All entries are stipulated economic net values at one common horizon. The
funding-only effect is $20$, the servicing-only effect is $15$, and their
joint effect is $5$. In general the decomposition is
\begin{align}
 E\{Y(1,1)-Y(0,0)\}
   &=E\{Y(1,0)-Y(0,0)\}+E\{Y(0,1)-Y(0,0)\}+I,\\
 I&=E\{Y(1,1)-Y(1,0)-Y(0,1)+Y(0,0)\}.
 \label{exec13:i3-interaction}
\end{align}
Here $I=-30$. Adding two correct isolated analyses and reporting $35$
would omit the interaction. The example does not assert that this particular
interaction occurs in a bank; it demonstrates information the proposed
portfolio consumer needs before combining answers.

Suppose the funding population has 100 accounts, the servicing population
60, and 20 accounts receive both actions. Assume the table applies to each
account, memberships are fixed before treatment, and accounts do not affect
one another. The total incremental value is
\[
 80(20)+40(15)+20(5)=2300,
\]
equivalently $100(20)+60(15)+20(-30)$. The naive separate total is
$2900$. Both evaluations and the equivalent decomposition are reproduced
exactly~\citep{exec13_i3_composition}. Origin identities are consequently
needed for more than deduplication: they establish which units receive the
joint policy. If linkage is uncertain, both membership and the relevant
interaction can be uncertain. If there are spillovers, the no-interference
assumption used by this sum must be replaced rather than hidden in the
aggregation code.

Before adding monetary outputs, the consumer must also identify whose cash
flow is being counted. An internal transfer of $3$ USD recorded as income by
a branch and cost by treasury sums to zero for the consolidated bank. The
branch target and bank target are different even though both are expressed in
USD. Similarly, a principal repayment may be relevant to liquidity while not
being income under the declared net-contribution measure. A state-space
quantity called a balance and a financial quantity called contribution cannot
be combined merely because both have currency units.

The inspected FinMarketing \texttt{CashEntry} distinguishes revenue, cost,
principal, liability, internal transfer and continuation, and retains owner,
cash date, economic basis and inclusion in value. Its \texttt{Value}
validation restricts owner and currency to the target, disallows principal,
liability and internal transfer from the net-contribution measure, and
requires continuation to start after the explicit cash horizon~\citep{exec13_i3_composition}. These are concrete source declarations for a
proposed valuation bridge. They do not establish all accounting policies or
cross-consumer correctness. CardNPV's continuation convention still requires
its own mapping, and a loan cash-flow valuation may legitimately include
principal under a different explicitly declared target.

A proposed combined workflow should therefore retain event and account
identity, economic owner, category, horizon and intervention membership before
aggregation. A manual reference can reconcile one small joint scenario across
the consumer records. An automated adapter should reproduce that result and
retain the same exceptions when an owner, horizon or intervention changes.
When a bridge cannot be justified, separate displayed values with a clear
non-additivity qualification are more informative than a fabricated total.
The resulting next experiment is specific: determine whether the consequential
joint response or population bridge is adequately supported, rather than
repeating isolated analyses whose individual accuracy cannot resolve it.


% SOURCE docs/release/2026-10-29/source/docs/writing-repair/2026-10-29/src/67dc0510606153fb.tex CHAR 0
\section{Standalone analysis and the three transport surfaces}

The same target restrictions apply when no sibling repository is involved.
A notebook, a batch script and a compatible MCP client are three ways of
submitting an analysis; none supplies missing causal premises. The library
and CLI chapter gives a complete supported request and its observed output,
including conditional identification and set-valued sensitivity results.
Its thirteen generated comparisons and associated lifecycle checks establish
a bounded local reference. Earlier MCP trials supply evidence for the
specific inspected local protocol and environment, rather than universal
compatibility with every MCP client or HTTP deployment~\citep{exec13_i2_examples,v12_in_runtime}. The separate later parity study in Section~\ref{exec16:semantic-parity} extends the local transport evidence to its own thirteen scientific requests and three lifecycle calls; it still does not demonstrate any of the consumer mappings proposed here.

For notebook use, the analyst should construct the target and scientific
context explicitly, save the full request and response, and derive any
displayed scalar from the response's declared type. The funding scenario
above is naturally a small local calculation: all cash assumptions and
coefficient interpretations are visible. A notebook cell that keeps only
the midpoint of a value range would make its positive result appear more
certain than the underlying information supports. If a graph analysis
returns the two-model effect set, a chart should preserve the two alternatives
or explicitly label an outer interval, rather than reporting their average
as an identified estimate.

For batch use, a manifest should give one immutable request identity per
business question and retain process exit status separately from the
scientific disposition. The common CLI accepts the operation and a versioned
JSON request; the complete executable listing appears in the preceding
chapter. A batch coordinator should distinguish malformed input, unavailable
dependencies, supported refusal and a computed result. A nonzero process
status does not imply a scientifically zero effect; a zero process status
does not imply identified benefit. The aggregate report must keep counts of
all submitted requests, including those with unavailable results. Reporting
average value only over successful calls changes the population whenever
success depends on difficult cases.

For MCP use, the client first negotiates the actual protocol and discovers
the available operations. It should then submit the same bounded scientific
request it would submit locally, keeping human choices about targets and
premises separate from tool mechanics. A tool can return a well-formed
unsupported or refused result. An agent must preserve that result and the
reason, rather than changing the sampling law or omitting the disputed
premise to obtain a number. The proposed graph and policy workflows may
need operations outside the current local implementation. Discovery of a
tool name is not evidence that it supports those domains.

The operating store is another boundary. Artifact handles are references
under a declared local custody root, not portable permissions to read any
file with the same name. A consumer-side export should preserve the content
digest, source role and permission context while mapping custody explicitly.
The proposed client should not pass an arbitrary sibling path as if every
service could access it. For a service wrapper, ownership, authorization,
schema negotiation and request/result lifecycle must be designed for the
actual deployment. Existing local stdio evidence does not establish network
isolation, multi-tenant behavior or a bank service level.

\section{A concrete conformance trial for each journey}

The proposed first trials start with the consumer's request and finish with
the object it displays. They compare both endpoints against an independently
specified scientific reference. This is stronger than testing whether an
encoder can decode its own output: the same mistaken unit convention can
survive a perfect round trip. Conversely, byte-for-byte equality is not always
required after a legitimate conversion. A dollars-to-millions exchange should
be equal after its declared inverse mapping, with a justified numerical
tolerance and the same target identity.

Table~\ref{exec13:i3-conformance} distinguishes the required reference from
the adapter behavior to be demonstrated. The calculations in this chapter
already supply finite examples for the scientific side. The table remains a
proposed consumer trial, since no sibling consumer was executed by those
calculations. The tests require actual exported requests, supported version
identities and review of the final presentation before an integration claim.

\begin{longtable}{p{.17\linewidth}p{.37\linewidth}p{.37\linewidth}}
\caption{Proposed consumer conformance cases.}
\label{exec13:i3-conformance}\\
\toprule Journey & Reference and successful exchange & Adverse exchange and required disposition \\
\midrule\endfirsthead
\toprule Journey & Reference and successful exchange & Adverse exchange and required disposition \\
\midrule\endhead
Funding & Four dated cash flows reproduce
Equation~\eqref{exec13:i3-funding}; dollars and millions give the same
mapped value. & Change basis-point convention, accrual period or exposure
response. Transform only with a justified identity; otherwise require a
new target or refuse. \\
Servicing & Two decision histories retain both policies, action costs and
terminal horizon; finite-law values are 29 and 43.6. & Make $L$ available
after $A_2$, substitute a smoothed state, or replace the comparator. Preserve
the original target and return unsupported information or a new analysis. \\
Discovery & Retain both specified models, effects $0$ and $0.6$, and net
values $-20$ and $40$. & A consumer requiring one scalar must receive an
explicit ambiguity disposition; no silent graph choice, mean effect or
favorable-value selection. \\
Observations & Preserve origin, inclusion, provenance and population;
enumeration yields expectation 50 for the expansion mean. & Set an inclusion
probability to zero, lose a shared map identity, or change owner. Require
additional information, a different target or a scoped refusal. \\
Standalone & Library, CLI and compatible MCP preserve the same supported
request and full scientific result. & Corrupt an artifact, change context,
cancel or retry. Distinguish invocation error, cancellation, replay and
scientific refusal without fabricating a result. \\
\bottomrule
\end{longtable}

Trials should also examine composite failure. A correct feature clock does
not rescue an incorrect population map; a correct graph representation does
not repair missing selection information. The acceptance object is the entire
declared workflow. Nevertheless, a failure should be localized: an unavailable
Java dependency should not prevent a reviewed Python file exchange, and an
unsupported longitudinal estimator should not block the exact funding-value
transformation. The common contract allows each consumer to see which part
is available, which scientific claim remains conditional and which next
input could change the answer.

For deterministic conversions, exact arithmetic or a stated floating-point
tolerance is appropriate. For a stochastic method, compare targets and
matched distributional behavior under a registered design, not a single
numerically similar run. Calibration, failure frequency, useful precision,
runtime and memory are distinct measurements. A faster adapter does not
repair undercoverage; a conservative whole-line set does not satisfy a
decision precision requirement simply because its coverage is high. The
operating chapter develops these resource and failure distinctions further.

\section{Choosing the first connection and the next experiment}

For this proposal, a reviewed manual artifact exchange is the baseline for
each selected consumer. It forces the target, units, clock and provenance to
be written down before automation hides them. The first automated connection
should address a frequent, stable exchange whose scientific interpretation
is already defensible. FinMarketing's source and target declarations and
dataml's artifact-oriented CLI provide plausible observation interfaces;
MacroFinance with the selected BayesFilter authority provides a plausible
funding route. Their priority follows from the banking questions, not from
an assertion that the mappings are complete.

A planning comparison can make the choice explicit. If a manual process
costs $c_m$ per admitted request, an automated route costs $c_a$, and design,
verification and maintenance have present cost $F$, then the cost advantage
over $n$ otherwise equivalent requests is
\[
  n(c_m-c_a)-F.
\]
The expression requires equivalent output quality, exception handling and
retention; it is not a measured return on investment. If $c_m\leq c_a$,
volume alone does not pay for the automation. If an adapter reduces meaningful
translation errors, their expected loss can be added only with an explicit
error model, baseline and uncertainty. Runtime for one arithmetic call omits
human specification, source review, data transfer, failures and revision cost.

The immediate informative experiment is therefore a small matched workflow
trial on public or clearly generated requests. Measure the manual and
automated routes against the same independently declared targets, preserve
every refusal and failed request, and include the adverse cases in
Table~\ref{exec13:i3-conformance}. Record where the actual consumer changes
or drops a field. This can justify building a broader adapter without the
cost of an authentic bank study. A later bank experiment is warranted when
the unresolved question concerns representative assignment, ownership,
precision, operating fit or value that the public trial cannot establish.

The reusable contribution is the target-preserving workflow, including its
failure states. The exact examples are reproducible with
\path{tools/workspace_examples_20261013.py}, using a fresh output path; they
require only the Python standard library and execute no sibling code~\citep{exec13_i3_examples}. They demonstrate why the proposed interfaces need
the information specified here. Actual adapter conformance, consumer
acceptance and bank applicability require their own evidence, attached to
the selected implementation and the claim being made.


% SOURCE docs/release/2026-10-29/source/docs/writing-repair/2026-10-29/src/d97a1b72be46db93.tex CHAR 928

\label{exec13:i3-end}

